% Generated draft: requires mathematical and source review.
\documentclass[11pt]{article}
\usepackage[margin=25mm]{geometry}
\usepackage{fontspec}
\setmainfont{Latin Modern Roman}
\usepackage{amsmath,amssymb,mathtools,unicode-math}
\setmathfont{Latin Modern Math}
\usepackage{enumitem,needspace,xurl,hyperref}
\setlength{\parindent}{0pt}
\setlength{\parskip}{5pt}
\setlength{\emergencystretch}{3em}
\newcommand{\N}{\mathbb{N}}
\newcommand{\Z}{\mathbb{Z}}
\newcommand{\Q}{\mathbb{Q}}
\newcommand{\R}{\mathbb{R}}
\newcommand{\C}{\mathbb{C}}
\newcommand{\F}{\mathbb{F}}
\newcommand{\A}{\mathbb{A}}
\newcommand{\PP}{\mathbb{P}}
\newcommand{\E}{\mathbb{E}}
\newcommand{\eps}{\varepsilon}
\newcommand{\dd}{\,\mathrm{d}}
\begin{document}
\section*{UnsolvedMath ID: 11000123}
\textbf{Source problem code:} catalogue entry
\subsection{Definitions and mathematical statement}
Let \(S_g\) be the orientable closed genus-\(g\) surface and \(N\) the closed nonorientable surface of genus \(g\), as in the accompanying source discussion. In the usual conventions \(\operatorname{Mod}_g\) uses orientation-preserving mapping classes of \(S_g\), \(\operatorname{Mod}(N)\) uses mapping classes of \(N\), and an asterisk denotes the extended orientable mapping class group, allowing orientation reversal. The allowed genus range and any exceptional conventions are not stated.

\emph{Problem 3.5.} Study homomorphisms in both directions,
\[
 \operatorname{Mod}_g\longrightarrow\operatorname{Mod}(N),
 \qquad\operatorname{Mod}(N)\longrightarrow\operatorname{Mod}_g.
\]
The source does not specify a classification up to conjugacy, a finite-image criterion, injective homomorphisms only, or another narrower target; none is silently imposed.

The source credits Birman and Chillingworth [3] with relating \(\operatorname{Mod}(N)\) to a subgroup of \(\operatorname{Mod}_{g-1}^*\), modulo an element of order two, through the orientation double cover \(S_{g-1}\to N\). It then states that certain finite-index subgroups of \(\operatorname{Mod}(N)\) are isomorphic to subgroups of \(\operatorname{Mod}_{g-1}\). These are supplied structural context. The exact lift subgroup and quotient in the first statement are not defined: the deck involution must be normal in the subgroup being quotiented, and a quotient of the entire extended mapping class group by an arbitrary order-two subgroup is not thereby defined. The draft does not replace the finite-index embedding statement by an embedding of the full \(\operatorname{Mod}(N)\), or confuse the cover's genus \(g-1\) with the genus \(g\) groups in the two requested directions.
\subsection{Short English statement}
Investigate homomorphisms between the orientable and nonorientable mapping class groups in both directions. The orientation-double-cover comparison supplies context involving genus \(g-1\) and finite-index subgroups, but does not itself classify those homomorphisms.
\subsection{Sources}
\begin{itemize}
\item Local UnsolvedMath registry, permanent ID 11000123, source code catalogue entry, Problem 3.5, PDF page 98. The chapter header identifies M. Korkmaz; no individual proposer field is supplied.
\item Farb, \emph{Problems on Mapping Class Groups and Related Topics}. Supplied URL: \url{https://www.math.uchicago.edu/~farb/papers/mcgbook.pdf}; recorded access 2026-07-29; extraction pdf-font-structure. The volume contains 329 explicitly typeset headings across chapters, with repeated local numbering disambiguated by page.
\item Birman and Chillingworth [3] is cited for the orientation-double-cover comparison. The full bibliographic entry, lift subgroup, exact quotient, and low-genus hypotheses are not provided.
\item Supplied status open and OPEN-TRIAGE. The summary mentions studied rigidity/finiteness results and says the problem is not fully settled, without precise references. Current primary-literature status and redistribution terms for the publicly accessible author/publisher copy remain NEEDS\_REVIEW.
\end{itemize}
\end{document}
